\documentclass[letterpaper]{article}
\usepackage[submission]{aaai2027}
\usepackage[hyphens]{url}
\usepackage{graphicx}
\urlstyle{rm}
\def\UrlFont{\rm}
\usepackage{natbib}
\usepackage{caption}
\frenchspacing
\usepackage{amsmath,amssymb}
\usepackage{booktabs}

\pdfinfo{
/TemplateVersion (2027.1)
}

\graphicspath{
  {../../../CoFiTok-internal/artifacts/figures/method_overview_2026-07-08/}
  {../../../CoFiTok-internal/artifacts/figures/visual_panels_2026-07-08/}
  {../../../CoFiTok-internal/artifacts/figures/final_visual_claim_2026-07-11/}
  {../../../CoFiTok-internal/artifacts/figures/final_visual_claim_2026-07-10/}
  {../../../CoFiTok-internal/artifacts/figures/ablation_hyperparameters_2026-07-11/}
  {../../../CoFiTok-internal/artifacts/figures/summary_2026-07-08/}
}

\newcommand{\method}{CoFiTok}
\newcommand{\epshat}{\hat{\epsilon}}
\newcommand{\xzero}{x_0}
\newcommand{\xt}{x_t}
\newcommand{\zk}{z^{(k)}}
\newcommand{\sk}{S_k}

\title{\method: Coarse-to-Fine Denoising Tokens for Pixel-Space Diffusion}
\author{Anonymous Submission}
\affiliations{}

\begin{document}
\maketitle

\begin{abstract}
Pixel-space diffusion models usually predict a full dense noise field with a
single monolithic network output. We introduce \method, a denoising-token
factorization that represents the dense noise prediction as an ordered sum of
negative-noise components. Each token is expanded by a restricted,
condition-free synthesis operator that only sees the current token, so prefix
sums form valid partial noise predictions and produce controllable partial
denoising results in pixel space. Across eight current datasets spanning
CIFAR-10, ImageNet-family 64x64 and 256x256 settings, FFHQ-64, and AFHQv2-64,
\method{} improves cross-batch prefix path AUC over an endpoint-only factorized control
on all datasets and avoids zero-token leakage that appears when the synthesis
operator is made deeper. In a two-seed ImageNet-256 confirmatory study, the
learned order ranks first among all 24 K4 permutations for both seeds while
preserving the endpoint sum. Frechet-style generation metrics favor EDM or direct dense
epsilon on most rows, so the supported claim is not unconditional generation
superiority; it is that ordered restricted dense-noise factorization yields
prefix-controllable denoising without relying on a VAE-style decoder.
\end{abstract}

\section{Introduction}

Diffusion models reverse a noising process by repeatedly predicting a dense
noise field. This dense prediction is effective, but it is also monolithic:
there is no direct way to ask how much of the denoising update has been
explained, whether the update has an interpretable coarse-to-fine order, or how
to stop after a valid partial prediction. Ordered visual tokenizers and
diffusion-autoregressive methods suggest that generation can benefit from
structured intermediate representations. The question in this paper is
narrower: can the dense pixel-space noise prediction itself be factorized into
ordered denoising tokens?

\method{} rewrites the predicted noise field as
\begin{equation}
  \epshat^{(K)} = \sum_{k=1}^{K} S_k\left(z^{(k)}\right),
  \label{eq:full-factorization}
\end{equation}
where each $z^{(k)}$ is a denoising token and $S_k$ expands only
that token into one dense component. Prefixes are meaningful by construction:
\begin{align}
  \epshat^{(1:m)} &= \sum_{k=1}^{m} S_k\left(z^{(k)}\right), \\
  \hat{x}^{(m)}_0 &= \frac{\xt - \sigma_t \epshat^{(1:m)}}{\alpha_t}.
  \label{eq:prefix-denoise}
\end{align}

The central design constraint is that the synthesis operator must not become a
decoder. In the default \method{} model, $S_k$ is token-only, bias-free,
shallow, local, and condition-free. It does not receive the noised image,
timestep, class label, text prompt, previous tokens, skip features, or a learned
constant. This restriction forces the token itself to carry the sample-specific
negative-noise information and gives the zero-token test a concrete meaning:
$S_k(0)=0$ by construction.

The current evidence supports a scoped method claim. \method{} learns ordered
denoising components across eight datasets, including ImageNet-family 64x64 and
256x256 settings and non-classification FFHQ/AFHQv2 rows. Compared with an
endpoint-only factorized predictor, \method{} consistently improves
  prefix/path alignment, while matched objective, prediction, and synthesis
  ablations expose why ordering and restricted synthesis matter. Generated-sample
distribution metrics do not show a stable unconditional generation-quality win,
which clarifies the present contribution: \method{} makes dense noise
prediction prefix-controllable without relying on a final VAE-style decoder.

\paragraph{Contributions.}
\begin{enumerate}
  \item We formulate dense pixel-space diffusion noise prediction as an ordered
  sum of denoising-token components.
  \item We introduce a restricted, condition-free synthesis operator whose
  token-only interface and bias-free construction prevent decoder-style
  shortcuts and make $S_k(0)=0$ testable.
  \item We train \method{} with a denoise-path objective that makes token
  prefixes approximate a partial denoising trajectory rather than only the final
  clean image.
  \item We evaluate prefix controllability, component ordering,
  non-degeneration, endpoint quality, matched short-run generation metrics, and
  tokenizer-related reconstruction baselines across the current eight-dataset
  evidence matrix.
\end{enumerate}

\section{Related Work}

\paragraph{Diffusion models and latent diffusion.}
Denoising diffusion probabilistic models learn to reverse a Gaussian noising
process by predicting noise targets \citep{ho2020denoising}. Latent diffusion
models move this process into a learned latent space and decode the final latent
through a VAE-style decoder \citep{rombach2022latentdiffusion}. \method{}
instead stays in the pixel-space denoising algebra and factorizes the dense
noise prediction itself. The \method{} token is not a final image latent; it is a
negative-noise component that directly contributes to the pixel-space
denoising update.

\paragraph{Trajectory ordering and latent-pixel interaction.}
Latent Forcing studies how changing the schedule between latent and pixel
variables changes the information revealed by the diffusion trajectory
\citep{baade2026latentforcing}. \method{} is inspired by this view of trajectory
ordering, but tokenizes the noise prediction inside the reverse trajectory
rather than introducing a separate latent-pixel trajectory.

\paragraph{Ordered visual tokenizers.}
Spectral Image Tokenizer, FlexTok, Matryoshka Representation Learning, M3, and
related nested-token approaches organize image or multimodal representations so
that prefixes or subsets can be useful
\citep{esteves2025spectral,bachmann2025flextok,kusupati2024matryoshka,cai2024matryoshkamultimodal}.
These methods operate on image representations or visual-token budgets.
\method{} operates on dense diffusion noise prediction. Its prefixes are valid
pixel-space denoising updates, not partial decodes from a general image
tokenizer.

\paragraph{Diffusion-autoregressive visual tokens.}
Selftok and D-AR are the closest conceptual neighbors because they connect
diffusion processes and autoregressive tokens \citep{wang2025selftok,gao2025dar}.
\method{} differs in the represented object and in the anti-degeneration
mechanism. \method{} tokens are continuous negative-noise components,
and the synthesis operator is deliberately weak and condition-free. The main
claim is not that diffusion steps are visual tokens, but that dense noise
prediction can be decomposed into ordered components whose prefixes are directly
usable in the diffusion denoising formula.

\section{Method}

\subsection{Forward Noising and Dense Prediction}

For a clean image $x_0$, timestep $t$, and Gaussian noise $\epsilon$, the forward
process is
\begin{equation}
  x_t = \alpha_t x_0 + \sigma_t \epsilon.
  \label{eq:forward}
\end{equation}
A conventional diffusion model predicts $\epsilon_\theta(x_t,t,c)$.
\method{} instead predicts the ordered component sum
in Eq.~\ref{eq:full-factorization}. The current implementation uses
unconditional models for the main MVP evidence. Class labels may exist in
datasets, but the restricted synthesis operator never receives labels, prompts,
timesteps, previous tokens, or image features. Figure~\ref{fig:method} shows the
dataflow and marks the blocked paths into $S_k$ that define the default method
boundary.

\begin{figure}[t]
  \centering
  \includegraphics[width=\linewidth]{method_overview.png}
  \caption{%
  \method{} factorizes dense pixel-space diffusion noise prediction into ordered
  denoising tokens. The predictor may be expressive, but each
  synthesis operator $S_k$ receives only the current token and is restricted to
  shallow, condition-free, bias-free local synthesis. Any prefix of components
  forms a valid partial noise prediction and therefore a valid pixel-space
  denoising result. Dashed blocked paths mark information that must not enter
  $S_k$, including $x_t$, timesteps, labels, prompts, previous tokens, skip
  features, and learned constants.}
  \label{fig:method}
\end{figure}

\subsection{Next Denoising Tokens}

The token predictor $T_k$ may be expressive. Our default implementation sees
the noised image, timestep, and previous tokens:
\begin{align}
  z^{(1)} &= T_1(x_t,t), \\
  z^{(k)} &= T_k(x_t,z^{(<k)},t), \quad k>1.
\end{align}
In the tiny MVP model, $T_k$ is a small U-Net-like predictor with token heads.
Autoregressive feedback is not part of the synthesis restriction or the core
factorization claim: a simultaneous variant may instead use
$z^{(k)}=T_k(x_t,t)$. We evaluate this choice directly.

\subsection{Restricted Synthesis Operator}

Each token is expanded independently:
\begin{equation}
  \epsilon^{(k)} = S_k\left(z^{(k)}\right).
\end{equation}
The default $S_k$ is a bias-free projection followed by a bias-free local
convolution. It is token-only and condition-free. Because no layer has bias and
there is no learned constant input, zero tokens map to zero components:
\begin{equation}
  S_k(0)=0.
\end{equation}
This restriction separates \method{} from a learned decoder. A deep-$S_k$
ablation is included only to test degeneration risk: it uses biased nonlinear
convolutions while preserving the token-only interface. This ablation can
improve apparent endpoint metrics but fails the zero-token diagnostic, so it is
not the default method.

\subsection{Prefix Denoising}

For any prefix length $m$, \method{} forms the prefix prediction and pixel-space
denoising result in Eq.~\ref{eq:prefix-denoise}. Thus every prefix is a valid
partial noise prediction and every prefix produces a pixel-space denoising
result. This is the primary interface exposed by the method.

\subsection{Denoise-Path Objective}

The current MVP objective is
\begin{equation}
  \mathcal{L}
  =
  \mathcal{L}_{\epsilon}
  + 0.15\,\mathcal{L}_{\mathrm{path\ prefix}}
  + 0.30\,\mathcal{L}_{\mathrm{path\ component}}.
\end{equation}
The denoise path interpolates from the zero-update $x_t$ reconstruction toward
coarse-to-fine spatial targets. The progress schedule uses power 1.5 in the
current best MVP configuration. This objective deliberately avoids forcing early
prefixes to be immediately clean; instead, prefixes are trained to follow a
partial denoising path.

\section{Experiments}

\subsection{Setup}

\paragraph{Datasets and models.}
We evaluate on CIFAR-10, Tiny ImageNet-200,
ImageNet-1K 64x64 HF, strict-source downsampled ImageNet-64, FFHQ-64,
AFHQv2-64, ImageNet-256 10\%, and full ImageNet-256. The HF ImageNet-64 row
and the strict-source downsampled ImageNet-64 row are reported separately; older
HF-derived results are not silently renamed as strict-source results.
All matched experiments use small unconditional models and short training
budgets. The primary model is K8 light \method{} at 64px and K4 at 256px. We
compare it with a parameter-matched direct dense predictor, an endpoint-only
factorized control, explicit objective/prediction/synthesis ablations, Improved
DDPM, and EDM. FlexTok/TiTok reconstruction and official D-AR/MAR/ReTok 50K
results are protocol-separated secondary context.

\paragraph{Protocols.}
The P0 rows match dataset, resolution, split, and nominal optimizer steps;
batch size, nominal images seen, parameter count, and NFE are reported because
the external methods are not compute matched. The direct dense control is the
strictest factorization control: it is parameter matched and uses the same
training budget, but is not claimed to be FLOP or latency matched.
The broad matrix uses locked per-dataset evaluation sizes. Repeated and
ablation claims use seeds 103/139 and 1,024 validation images at fixed $t=500$.
All power-sweep models are evaluated against the same progress-$1.5$ path.

\paragraph{Metrics and evidence tiers.}
We report endpoint PSNR and MSE, path AUC, effective token count, and zero-input
energy. Lower path AUC means better alignment of all prefixes
to the prescribed partial-denoising path; endpoint metrics alone do not measure
that interface. Frechet-style values are internal comparison metrics rather
than SOTA FID claims. The locked evidence matrix contains 80 matched P0 cells,
16 tokenizer eval-only cells, and 24 explicitly protocol-blocked cross-task
cells; full settings and provenance are in the supplement.

\subsection{Results}

\paragraph{All-dataset evidence.}

Table~\ref{tab:cofitok-diagnostics-current} isolates the supported claim. C is
\method{}, D is the parameter-matched direct dense predictor, and E is the
endpoint-only factorized control. \method{} has lower path AUC than E on all
eight datasets and exactly preserves the restricted zero-token contract. Its
endpoint PSNR is close to, but not uniformly better than, direct dense.

\begin{table*}[t]
  \centering
  \scriptsize
  \caption{Diagnostics for the supported \method{} claim. Higher PSNR is better;
  lower path AUC and zero ratio are better.}
  \label{tab:cofitok-diagnostics-current}
  \resizebox{\textwidth}{!}{%
  \begin{tabular}{lrrrrrrrr}
    \toprule
    Dataset & C PSNR & D PSNR & C path & E path & C effK & E effK & C zero & deep-$S_k$ zero \\
    \midrule
    \input{../../../CoFiTok-internal/artifacts/reports/aaai27_experiment_tables_2026-07-11/all_dataset_result_rows.tex}\\
    \bottomrule
  \end{tabular}%
  }
\end{table*}

\begin{figure*}[t]
  \centering
  \begin{minipage}[t]{0.49\textwidth}
    \centering
    \includegraphics[width=\linewidth]{figure2a_tiny_prefix.png}
    \par\smallskip
    {\footnotesize\textbf{(a)} Tiny-IN: endpoint-only (left), \method{} (right).}
  \end{minipage}\hfill
  \begin{minipage}[t]{0.49\textwidth}
    \centering
    \includegraphics[width=\linewidth]{figure2b_imagenet64_prefix.png}
    \par\smallskip
    {\footnotesize\textbf{(b)} ImageNet-64: endpoint-only (left), \method{} (right).}
  \end{minipage}

  \medskip
  \begin{minipage}[t]{0.49\textwidth}
    \centering
    \includegraphics[width=\linewidth]{figure2c_restricted_synthesis.png}
    \par\smallskip
    {\footnotesize\textbf{(c)} Restricted $S_k$: ordered (left), shuffled (right).}
  \end{minipage}\hfill
  \begin{minipage}[t]{0.49\textwidth}
    \centering
    \includegraphics[width=\linewidth]{figure2d_deep_synthesis.png}
    \par\smallskip
    {\footnotesize\textbf{(d)} Deep $S_k$: ordered (left), shuffled (right).}
  \end{minipage}
  \caption{%
  Qualitative prefix and degeneration evidence assembled from four standalone
  subfigure files. In (a--b), each half shows the same examples across the
  prefix path; \method{} produces more organized intermediate denoising than
  endpoint-only factorization. In (c--d), shuffling breaks sample alignment,
  while deep synthesis can carry structure that the restricted, zero-preserving
  operator cannot. Quantitative path and zero-token results are reported in
  Tables~\ref{tab:cofitok-diagnostics-current} and~\ref{tab:component-ablation}.}
  \label{fig:final-visual}
\end{figure*}

\paragraph{Endpoint quality and prefixes.}

At the same 20k training budget, K8 light \method{} nearly matches the
endpoint-only factorized control in endpoint $x_0$ MSE at $t=500$ while greatly
improving path alignment.

\begin{table*}[t]
  \centering
  \caption{Matched 20k two-seed comparison (mean $\pm$ sample standard
  deviation over seeds 103 and 139). Endpoint $x_0$ MSE and path AUC use
  1,024 validation images at $t=500$; lower is better. Higher is better for
  effective token count.}
  \label{tab:fair20k}
  \small
  \begin{tabular}{llrrrr}
    \toprule
    Dataset & Method & Endpoint $x_0$ MSE $\downarrow$ & Path AUC $\downarrow$ & Eff. K $\uparrow$ & Zero ratio $\downarrow$ \\
    \midrule
    \input{../../../CoFiTok-internal/artifacts/reports/long_budget_repeat_2026-07-11/long_budget_repeat_rows.tex}\\
    \bottomrule
  \end{tabular}
\end{table*}

Across the four paired dataset-seed comparisons, \method{} lowers path AUC in
4/4 cases, with a mean relative reduction of 96.18\%. This structural gain has
a measurable endpoint tradeoff: endpoint $x_0$ MSE increases by 4.20\% on average and
is not lower in any of the four pairs. Figure~\ref{fig:final-visual} shows the
corresponding qualitative prefix grids.
In the full-validation deterministic sweep, K8 light also improves MSE AUC and
low-res Frechet proxy AUC on both Tiny ImageNet-200 and ImageNet-64 HF, while
the endpoint-only factorized control keeps the lower endpoint MSE.

\paragraph{Learned order and ImageNet-256 scaling.}

Order ablations keep the same component sum at the endpoint but change the
prefix trajectory. This makes order evaluation a direct test of whether prefixes
are organized rather than an endpoint-quality test.

Table~\ref{tab:imagenet256-confirmatory-order} gives the preregistered
ImageNet-256 K4 20k two-seed test. Besides a fixed random permutation and
reverse order, it exhaustively evaluates all 24 component permutations over
1,024 images and reports a paired-bootstrap confidence interval for the mean of
the 23 non-identity orders. Because every order sums the same components, the
endpoint is preserved; the test isolates the organization of the prefix path.

\begin{table*}[t]
  \centering
  \scriptsize
  \caption{ImageNet-256 K4 20k confirmatory order and endpoint results over
  1,024 validation images at $t=500$.}
  \label{tab:imagenet256-confirmatory-order}
  \resizebox{\textwidth}{!}{%
  \begin{tabular}{rrrrrrrrr}
    \toprule
    Seed & E path & C ordered & non-ID mean & reverse & rank/24 & $\Delta$ CI low & D endpoint & C endpoint \\
    \midrule
    \input{../../../CoFiTok-internal/artifacts/reports/aaai27_experiment_tables_2026-07-11/imagenet256_confirmatory_rows.tex}\\
    \bottomrule
  \end{tabular}%
  }
\end{table*}

The preregistered confirmatory test passes all seven gates. For both seeds, the
learned order ranks 1st of 24, every non-identity permutation has worse path
AUC, and the paired-bootstrap lower bounds for non-identity minus ordered AUC
are positive (0.2489 and 0.2400). The endpoint component sum is invariant up to
$3\times10^{-15}$ MSE. Relative to the parameter-matched direct dense control,
\method{} also lowers endpoint $x_0$ MSE by 6.73\% on average in this
ImageNet-256 confirmatory setting.

\paragraph{Generation and related-method context.}
Cross-batch validation confirms that endpoint-only factorization remains a
strong endpoint baseline. Generated-sample metrics do not establish stable
unconditional superiority: the matched controls favor endpoint-only on Tiny
ImageNet-200 and ImageNet-64 HF under the internal streamed protocol, and EDM
or direct dense wins the broad Frechet-style rows. Official D-AR/MAR/ReTok 50K
and FlexTok/TiTok reconstruction results remain secondary protocol tiers in the
supplement; none is used as a matched \method{} superiority claim.

\subsection{Ablations}

\paragraph{Matched component protocol.}
We isolate the K8 method components on ImageNet-64 HF with identical 5k-step
training, seeds 103/139, and 1,024-image evaluation at $t=500$. The full model
uses $(\lambda_p,\lambda_c,p)=(0.15,0.30,1.5)$. Endpoint-only removes both path
terms; the next two rows remove one term at a time; simultaneous prediction
removes token feedback; deep synthesis replaces restricted $S_k$ with a biased
nonlinear decoder. Table~\ref{tab:component-ablation} reports mean $\pm$ sample
standard deviation.

\begin{table*}[t]
  \centering
  \small
  \caption{Matched two-seed component ablations on ImageNet-64 HF. Lower is
  better for endpoint MSE, path AUC, and zero ratio; higher effective K indicates
  broader component use. All path values use the same progress-$1.5$ target.}
  \label{tab:component-ablation}
  \begin{tabular}{lrrrr}
    \toprule
    Variant & Endpoint $x_0$ MSE $\downarrow$ & Path AUC $\downarrow$ & Eff. K $\uparrow$ & Zero ratio $\downarrow$ \\
    \midrule
    \input{../../../CoFiTok-internal/artifacts/reports/aaai27_ablation_2026-07-11/component_ablation_rows.tex}\\
    \bottomrule
  \end{tabular}
\end{table*}

\paragraph{Objective, prediction, and synthesis components.}
Relative to endpoint-only factorization, the full objective lowers mean path AUC
by 93.0\% with a 2.0\% endpoint-MSE cost. Relative to removing only the prefix
or component path term, it lowers path AUC by 45.7\% and 37.9\%, respectively,
while endpoint MSE changes by less than 0.2\%. Both path terms therefore improve
the claimed interface rather than endpoint fit. Sequential feedback is a modest
implementation gain, not a core requirement: full improves the two-seed mean by
2.8\% in path AUC and 1.0\% in endpoint MSE over simultaneous heads, but wins
path AUC on only one seed. Deep synthesis attains lower path AUC yet has zero
ratio $0.0485\pm0.0050$ versus exactly zero for restricted $S_k$; its apparent
gain violates the identifiability contract. Figure~\ref{fig:final-visual}
visualizes zero/shuffle/deep-$S_k$ diagnostics.

\paragraph{Hyperparameter sensitivity.}
Figure~\ref{fig:ablation-hyperparameters} varies the two path-loss weights,
progress power, and token count. Every point is a two-seed mean with sample
standard-deviation bars; the dashed line marks the default. Crucially, the power
sweep is evaluated against one common progress-$1.5$ path rather than allowing
each training target to define its own metric. Raising either path weight
steadily improves path AUC with little endpoint movement; $(0.15,0.30)$ is a
moderate tradeoff rather than the path-only optimum. Power 1.5 is the clear
minimum under the common target. K8 gives the lowest mean path AUC, while K4 has
slightly lower endpoint MSE and K16 adds cost without a 5k-step path gain.

\begin{figure*}[t]
  \centering
  \includegraphics[width=\textwidth]{aaai27_ablation_hyperparameters.pdf}
  \caption{Matched ImageNet-64 HF hyperparameter ablations. Top: denoise-path
  AUC; bottom: endpoint $x_0$ MSE. Curves vary prefix weight $\lambda_p$,
  component weight $\lambda_c$, progress power $p$, and token count $K$.
  Lower is better; bars show sample standard deviation over seeds 103/139.}
  \label{fig:ablation-hyperparameters}
\end{figure*}

\paragraph{Component-separation stress test.}

We treat component decorrelation as a pressure-term ablation rather than a
default objective. The headline \method{} rows use the no-decorrelation light
denoise-path objective because it most directly measures ordered prefix
denoising. On Tiny ImageNet-200, a conservative schedule
(\texttt{weight=0.003}, start step 3000, warmup 1500, full weight at step 4500
of 5000) reduces component correlation on two seeds while keeping the
ordered-prefix penalty small.

Across the two seeds, scheduled decor 0.003 changes ordered path AUC by +0.0036,
mean absolute component cosine by -0.0328, final MSE by -0.0002, and effective
token count by +0.0336. Thus mild separation pressure can make components less
correlated without materially breaking ordered prefix behavior. Stronger decor
weights reduce cosine further but impose larger path-AUC penalties, so we keep
them as stress-test ablations rather than the default \method{} objective.

\section{Discussion}

\method{} suggests that dense pixel-space diffusion noise prediction can be
treated as a sequence of ordered denoising-token components. The restricted
synthesis operator is central: it gives the prefix interface meaning and
prevents the method from becoming a generic latent decoder. The denoise-path
objective is also central because early prefixes should represent partial
denoising rather than clean reconstructions.

The strongest empirical result is the all-dataset prefix-control diagnostic:
\method{} has lower cross-batch path AUC than the endpoint-only factorized control on all
eight datasets, while matched objective, prediction, order, and deep-$S_k$
ablations identify the roles of path supervision and restricted synthesis. The most important negative result is that
generated-sample Frechet-style metrics do not show a stable unconditional
generation-quality win. Together, these results define the current scope of the
contribution.

\section{Limitations}

The current models are small and trained with short MVP budgets. Frechet-style
sample-quality results are mixed and do not show a robust \method{} generation
win; EDM is stronger on most generation rows, while direct dense prediction
provides the strict factorization control. P1 FlexTok and
TiTok rows use official pretrained reconstruction at 256x256 and are not
retrained generation baselines. D-AR, MAR, and ReTok use separate official
ImageNet-256 eval-only protocols and remain protocol-blocked under the
matched-dataset/step training definition. Seed repeats cover the main
endpoint-only/K8 comparison but not every ablation. The AAAI-27 adaptation is an
anonymous submission build and still needs author metadata and final checklist
decisions before any camera-ready version. Sequential token prediction can also
cost more latency and memory than a direct dense head; the current audit reports
parameters, batch, nominal training images, and NFE, but not a standardized
cross-implementation FLOP or wall-clock benchmark.

\section{Conclusion}

\method{} factorizes pixel-space diffusion noise prediction into an ordered
  sequence of denoising tokens. Each token is expanded by a restricted,
condition-free synthesis operator into a dense negative-noise component, and
prefix sums yield valid partial denoising results. Across the current
eight-dataset evidence matrix, \method{} consistently improves prefix-control
diagnostics over an endpoint-only factorized control and passes the restricted zero-token
diagnostic where deep synthesis leaks token-independent structure. The current
evidence therefore supports a scoped contribution about ordered restricted
dense-noise factorization, while leaving raw unconditional generation quality
to stronger monolithic diffusion baselines. The next step is to scale the
protocol with stronger backbones, longer training, and broader seed coverage.

\bibliography{../../references}

\end{document}
